% Options for packages loaded elsewhere
\PassOptionsToPackage{unicode}{hyperref}
\PassOptionsToPackage{hyphens}{url}
\documentclass[
  a4paper,
]{ltjsarticle}
\usepackage{xcolor}
\usepackage[margin=25mm]{geometry}
\usepackage{amsmath,amssymb}
\setcounter{secnumdepth}{-\maxdimen} % remove section numbering
\usepackage{iftex}
\ifPDFTeX
  \usepackage[T1]{fontenc}
  \usepackage[utf8]{inputenc}
  \usepackage{textcomp} % provide euro and other symbols
\else % if luatex or xetex
  \usepackage{unicode-math} % this also loads fontspec
  \defaultfontfeatures{Scale=MatchLowercase}
  \defaultfontfeatures[\rmfamily]{Ligatures=TeX,Scale=1}
\fi
\usepackage{lmodern}
\ifPDFTeX\else
  % xetex/luatex font selection
\fi
% Use upquote if available, for straight quotes in verbatim environments
\IfFileExists{upquote.sty}{\usepackage{upquote}}{}
\IfFileExists{microtype.sty}{% use microtype if available
  \usepackage[]{microtype}
  \UseMicrotypeSet[protrusion]{basicmath} % disable protrusion for tt fonts
}{}
\makeatletter
\@ifundefined{KOMAClassName}{% if non-KOMA class
  \IfFileExists{parskip.sty}{%
    \usepackage{parskip}
  }{% else
    \setlength{\parindent}{0pt}
    \setlength{\parskip}{6pt plus 2pt minus 1pt}}
}{% if KOMA class
  \KOMAoptions{parskip=half}}
\makeatother
\usepackage{color}
\usepackage{fancyvrb}
\newcommand{\VerbBar}{|}
\newcommand{\VERB}{\Verb[commandchars=\\\{\}]}
\DefineVerbatimEnvironment{Highlighting}{Verbatim}{commandchars=\\\{\}}
% Add ',fontsize=\small' for more characters per line
\newenvironment{Shaded}{}{}
\newcommand{\AlertTok}[1]{\textcolor[rgb]{1.00,0.00,0.00}{\textbf{#1}}}
\newcommand{\AnnotationTok}[1]{\textcolor[rgb]{0.38,0.63,0.69}{\textbf{\textit{#1}}}}
\newcommand{\AttributeTok}[1]{\textcolor[rgb]{0.49,0.56,0.16}{#1}}
\newcommand{\BaseNTok}[1]{\textcolor[rgb]{0.25,0.63,0.44}{#1}}
\newcommand{\BuiltInTok}[1]{\textcolor[rgb]{0.00,0.50,0.00}{#1}}
\newcommand{\CharTok}[1]{\textcolor[rgb]{0.25,0.44,0.63}{#1}}
\newcommand{\CommentTok}[1]{\textcolor[rgb]{0.38,0.63,0.69}{\textit{#1}}}
\newcommand{\CommentVarTok}[1]{\textcolor[rgb]{0.38,0.63,0.69}{\textbf{\textit{#1}}}}
\newcommand{\ConstantTok}[1]{\textcolor[rgb]{0.53,0.00,0.00}{#1}}
\newcommand{\ControlFlowTok}[1]{\textcolor[rgb]{0.00,0.44,0.13}{\textbf{#1}}}
\newcommand{\DataTypeTok}[1]{\textcolor[rgb]{0.56,0.13,0.00}{#1}}
\newcommand{\DecValTok}[1]{\textcolor[rgb]{0.25,0.63,0.44}{#1}}
\newcommand{\DocumentationTok}[1]{\textcolor[rgb]{0.73,0.13,0.13}{\textit{#1}}}
\newcommand{\ErrorTok}[1]{\textcolor[rgb]{1.00,0.00,0.00}{\textbf{#1}}}
\newcommand{\ExtensionTok}[1]{#1}
\newcommand{\FloatTok}[1]{\textcolor[rgb]{0.25,0.63,0.44}{#1}}
\newcommand{\FunctionTok}[1]{\textcolor[rgb]{0.02,0.16,0.49}{#1}}
\newcommand{\ImportTok}[1]{\textcolor[rgb]{0.00,0.50,0.00}{\textbf{#1}}}
\newcommand{\InformationTok}[1]{\textcolor[rgb]{0.38,0.63,0.69}{\textbf{\textit{#1}}}}
\newcommand{\KeywordTok}[1]{\textcolor[rgb]{0.00,0.44,0.13}{\textbf{#1}}}
\newcommand{\NormalTok}[1]{#1}
\newcommand{\OperatorTok}[1]{\textcolor[rgb]{0.40,0.40,0.40}{#1}}
\newcommand{\OtherTok}[1]{\textcolor[rgb]{0.00,0.44,0.13}{#1}}
\newcommand{\PreprocessorTok}[1]{\textcolor[rgb]{0.74,0.48,0.00}{#1}}
\newcommand{\RegionMarkerTok}[1]{#1}
\newcommand{\SpecialCharTok}[1]{\textcolor[rgb]{0.25,0.44,0.63}{#1}}
\newcommand{\SpecialStringTok}[1]{\textcolor[rgb]{0.73,0.40,0.53}{#1}}
\newcommand{\StringTok}[1]{\textcolor[rgb]{0.25,0.44,0.63}{#1}}
\newcommand{\VariableTok}[1]{\textcolor[rgb]{0.10,0.09,0.49}{#1}}
\newcommand{\VerbatimStringTok}[1]{\textcolor[rgb]{0.25,0.44,0.63}{#1}}
\newcommand{\WarningTok}[1]{\textcolor[rgb]{0.38,0.63,0.69}{\textbf{\textit{#1}}}}
\usepackage{longtable,booktabs,array}
\newcounter{none} % for unnumbered tables
\usepackage{calc} % for calculating minipage widths
% Correct order of tables after \paragraph or \subparagraph
\usepackage{etoolbox}
\makeatletter
\patchcmd\longtable{\par}{\if@noskipsec\mbox{}\fi\par}{}{}
\makeatother
% Allow footnotes in longtable head/foot
\IfFileExists{footnotehyper.sty}{\usepackage{footnotehyper}}{\usepackage{footnote}}
\makesavenoteenv{longtable}
\setlength{\emergencystretch}{3em} % prevent overfull lines
\providecommand{\tightlist}{%
  \setlength{\itemsep}{0pt}\setlength{\parskip}{0pt}}
\usepackage{bookmark}
\IfFileExists{xurl.sty}{\usepackage{xurl}}{} % add URL line breaks if available
\urlstyle{same}
\hypersetup{
  pdftitle={Chapter 2: The Sweep --- Stage 1+2+3 Dynamics and the N=3..40 Runs (H⊥/H Denominator-Control Figure)},
  pdfauthor={Noriaki Kihara (WF System Co., Ltd.) ORCID 0009-0004-6753-4020},
  hidelinks,
  pdfcreator={LaTeX via pandoc}}

\title{Chapter 2: The Sweep --- Stage 1+2+3 Dynamics and the N=3..40
Runs (H⊥/H Denominator-Control Figure)}
\author{Noriaki Kihara (WF System Co., Ltd.) ORCID 0009-0004-6753-4020}
\date{September 5, 2026 Version DOI 10.5281/zenodo.22317636}

\begin{document}
\maketitle

(Chapter 2 of the reproduction papers for the N=3..40 stage-1+2+3 sweep.
Shares the Concept DOI 10.5281/zenodo.22317635 with the overview paper;
Version DOI 10.5281/zenodo.22317636. Equation numbers continue from
Chapter 1 (Eqs. 1--9). Code blocks are quoted verbatim from the
programs; comments inside them remain in Japanese, as in the originals.)

\subsection{1. Purpose}\label{purpose}

The central purpose of this series of numerical experiments is
\textbf{to clarify the mechanism by which inflation-like development
occurs --- the phenomenon in which a seed too small to be measured
(dormant fraction H⊥/H \textasciitilde{} 10\^{}-30) is exponentially
amplified by the dynamics alone over dozens of orders of magnitude until
saturation.} Since this phenomenon was observed in the July canonical
runs (N = 40, 300, 1000), the series has isolated, one factor at a time,
which components of the dynamics are necessary conditions for this
development (the confirmed result is the stage-1+2+3 composition; the
audit record is Appendix A).

As the culmination, this chapter takes the 38 static parents (N=3..40)
generated and verified in Chapter 1 as initial data, runs the
stage-1+2+3 dynamics (phase-only, imaginary-part-only generator; real
orthogonal rotation; fixed Δτ=2π/den) for N=3..40 × 6 denominators × 500
steps, records at every step the indicator of inflation-like
development, the \textbf{dormant fraction H⊥/H} (definition and its
justification as an indicator: §2.5), the total energy H\_total, and the
zero closure \textbar z\^{}Tz\textbar/H, and fixes how this development
appears across the whole range of N and the sweep of the clock
(denominator) in the target figure
\texttt{fig\_Hperp\_denominator\_controls\_with\_124\_N3\_N40\_stage123.png}.

The descriptive purpose of this paper is complete reproducibility and
the fixing of the correspondence among equations, programs, and data; it
contains no physical interpretation beyond the observed facts. The
grounds for the stage-1+2+3 composition are placed in \textbf{Appendix A
(audit table of deletion controls)} and \textbf{Appendix B (program
lineage)} so as not to interrupt the main text.

\subsection{2. Theoretical Background}\label{theoretical-background}

The dynamics of this chapter takes the dynamics of the July canonical
(old) program as its starting point and reconstructs it on a new
architecture, decomposed into three components: \textbf{stage 1, stage
2, and stage 3}. The theoretical background is therefore described in
the order (2.1) the common state space, (2.2) the mathematics of the old
program's dynamics, (2.3) the mathematics of this program's dynamics
(stage decomposition), (2.4) the mathematical differences between the
old and new rotation maps, and (2.5) measured quantities and
conservation laws.

\subsubsection{2.1 State Space and Adjacency Structure (common to old
and
new)}\label{state-space-and-adjacency-structure-common-to-old-and-new}

\textbf{(Eq. 10) Line-graph adjacency matrix} --- For edges e=(i,j),
f=(k,l) of the complete graph K\_N, A\_ef = 1 when e and f share a
vertex, 0 otherwise, 0 on the diagonal (M×M, real symmetric). The edge
ordering is the upper-triangular order identical to Eq. 1. The state is
z ∈ C\^{}M (M = N(N−1)/2); the edge phase is θ\_e = arg z\_e.

\subsubsection{2.2 Mathematics of the Old Program (July Canonical)
Dynamics}\label{mathematics-of-the-old-program-july-canonical-dynamics}

The old program is defined by the engine
\texttt{run\_n\_scaling\_lowrank\_v1.py} bundled in this package (a
bit-identical copy of the canonical). One step is the composition of the
following three parts.

\textbf{(Eq. 11) Old generator (phase-difference sine, real
antisymmetric)} --- On vertex-sharing edge pairs (e,f):

\begin{verbatim}
K_ef(θ) = cos θ_e sin θ_f − sin θ_e cos θ_f = sin(θ_f − θ_e)
\end{verbatim}

The generator is rebuilt at every step from the current phases θ = arg
z. The amplitudes \textbar z\_e\textbar{} do not enter. The vertex
decomposition K = C S\^{}T − S C\^{}T = W J W\^{}T (rank ≤ 2N) is the
same as Eq. 3 of Chapter 1.

\textbf{(Eq. 12) Power-iteration estimate of σ\_max (approximate,
history-dependent)} --- A real vector wp is carried across steps:

\begin{verbatim}
wp ← −K(K wp)/‖−K(K wp)‖   (3 iterations)
σ^_max = ‖K wp‖
\end{verbatim}

Since −K\^{}2 is positive semidefinite symmetric with largest eigenvalue
σ\_max\^{}2, σ\^{}\_max is a \textbf{power-iteration approximation} of
σ\_max. Only 3 iterations are used, and wp is inherited from the
previous step (warm start). Strictly speaking, the old dynamics is
therefore not a map of z alone but a \textbf{system with a hidden state
wp}, and σ\^{}\_max(t) is an approximate sequence depending on the
trajectory history.

Engine implementation (\texttt{run\_n\_scaling\_lowrank\_v1.py}):

\begin{Shaded}
\begin{Highlighting}[]
   \DecValTok{122}      \KeywordTok{def}\NormalTok{ sigma\_max\_power(}\VariableTok{self}\NormalTok{, wp, iters}\OperatorTok{=}\DecValTok{3}\NormalTok{):}
   \DecValTok{123}          \StringTok{"""warm{-}start 冪反復による σ\_max 推定。wp は前ステップのベクトル（更新して返す）。"""}
   \DecValTok{124}          \ControlFlowTok{for}\NormalTok{ \_ }\KeywordTok{in} \BuiltInTok{range}\NormalTok{(iters):}
   \DecValTok{125}\NormalTok{              y }\OperatorTok{=} \VariableTok{self}\NormalTok{.kmatvec(wp)}
   \DecValTok{126}\NormalTok{              wp }\OperatorTok{=} \OperatorTok{{-}}\VariableTok{self}\NormalTok{.kmatvec(y)}
   \DecValTok{127}\NormalTok{              nrm }\OperatorTok{=}\NormalTok{ np.linalg.norm(wp)}
   \DecValTok{128}              \ControlFlowTok{if}\NormalTok{ nrm }\OperatorTok{==} \FloatTok{0.0}\NormalTok{:}
   \DecValTok{129}                  \ControlFlowTok{return} \FloatTok{0.0}\NormalTok{, wp}
   \DecValTok{130}\NormalTok{              wp }\OperatorTok{=}\NormalTok{ wp }\OperatorTok{/}\NormalTok{ nrm}
   \DecValTok{131}\NormalTok{          sig }\OperatorTok{=}\NormalTok{ np.linalg.norm(}\VariableTok{self}\NormalTok{.kmatvec(wp))}
   \DecValTok{132}          \ControlFlowTok{return} \BuiltInTok{float}\NormalTok{(sig), wp}
\end{Highlighting}
\end{Shaded}

\textbf{(Eq. 13) Old one-step map (Cayley transform, σ-normalized
clock)} --- With the constant γ = tan(π/144):

\begin{verbatim}
K̃ = K / σ^_max
z ← (I − γ K̃)^-1 (I + γ K̃) z
\end{verbatim}

Engine implementation (2N×2N solve via Woodbury):

\begin{Shaded}
\begin{Highlighting}[]
   \DecValTok{134}      \KeywordTok{def}\NormalTok{ cayley\_step(}\VariableTok{self}\NormalTok{, z, sigma):}
   \DecValTok{135}          \StringTok{"""z ← (I{-}γK̃)\^{}\{{-}1\}(I+γK̃) z, K̃ = K/σ。Woodbury で O(N\^{}3)。"""}
   \DecValTok{136}\NormalTok{          gn }\OperatorTok{=}\NormalTok{ GAMMA }\OperatorTok{/}\NormalTok{ sigma}
   \DecValTok{137}\NormalTok{          r }\OperatorTok{=}\NormalTok{ z }\OperatorTok{+}\NormalTok{ gn }\OperatorTok{*} \VariableTok{self}\NormalTok{.kmatvec(z)}
   \DecValTok{138}\NormalTok{          A2 }\OperatorTok{=}\NormalTok{ (sigma }\OperatorTok{/}\NormalTok{ GAMMA) }\OperatorTok{*} \VariableTok{self}\NormalTok{.J }\OperatorTok{+} \VariableTok{self}\NormalTok{.G}
   \DecValTok{139}\NormalTok{          rhs }\OperatorTok{=} \VariableTok{self}\NormalTok{.wt(r)}
   \DecValTok{140}\NormalTok{          y }\OperatorTok{=}\NormalTok{ np.linalg.solve(A2, rhs)}
   \DecValTok{141}          \ControlFlowTok{return}\NormalTok{ r }\OperatorTok{{-}} \VariableTok{self}\NormalTok{.w(y)}
\end{Highlighting}
\end{Shaded}

\textbf{(Eq. 14) Old eigenplane rotation angle} --- The eigenvalues of K
are purely imaginary pairs ±iσ\_k (σ\_k ≥ 0). The Cayley factor for the
eigenvalues ±iσ\_k/σ\^{}\_max of K̃ is

\begin{verbatim}
(1 + iγσ_k/σ^_max) / (1 − iγσ_k/σ^_max)   (modulus 1)
\end{verbatim}

so the one-step rotation angle of eigenplane k is

\begin{verbatim}
φ_k = 2 arctan( γ · σ_k / σ^_max )
\end{verbatim}

The rotation angle of the fastest plane (σ\_k = σ\^{}\_max) is always
fixed at φ\_max = 2 arctan(tan(π/144)) = \textbf{π/72}. This is an
\textbf{adaptive (σ-normalized) clock} that ``ticks time by the system's
fastest mode'', and the spectral dependence of the rotation angle
undergoes \textbf{nonlinear compression} by arctan.

\textbf{(Eq. 15) Old conserved quantities} --- Since K is real
antisymmetric, the Cayley transform O = (I−γK̃)\^{}-1(I+γK̃) is a
\textbf{real orthogonal matrix} (O\^{}T O = I). A real orthogonal matrix
acts identically on the real and imaginary parts of z = x + iy, so

\begin{verbatim}
‖Oz‖^2 = ‖z‖^2          (norm conservation)
(Oz)^T(Oz) = z^T O^TO z = z^Tz   (exact conservation of the zero closure z^Tz)
\end{verbatim}

\textbf{(Eq. 15') Relative equilibrium (parent)} --- The parent v lies
on an eigenplane of K(arg v) (Chapter 1, Eqs. 4--6); on the flow dZ/dt =
KZ we have KZ = −iσZ, i.e., Z(t) = e\^{}\{−iσt\} Z(0). All edge phases
advance by the same amount, so phase differences are invariant, hence
K(θ(t)) = K(θ(0)). The parent is a relative equilibrium that ``rotates
rigidly without changing phase differences''.

\subsubsection{2.3 Mathematics of This Program's Dynamics (Definitions
of Stages 1, 2,
3)}\label{mathematics-of-this-programs-dynamics-definitions-of-stages-1-2-3}

The stage numbering follows the definitions of the N=40 single-factor
experiment series (Appendix A). \textbf{Stage 1 provides the baseline
architecture}, and \textbf{stages 2 and 3 are the two modifications
applied to it}.

\textbf{(Eq. 16) Stage 1: baseline architecture (explicit matrix,
spectral map, fixed Δτ clock)} --- The generator is constructed
explicitly as a Hermitian matrix H, and one step is given by the exact
spectral exponential map

\begin{verbatim}
H = V diag(w) V†,   z ← V e^{−i Δτ w} V† z
\end{verbatim}

The clock is a \textbf{fixed external clock}

\begin{verbatim}
Δτ = 2π/den,   den ∈ { N−2, N−1, N, N+1, N+2 } ∩ N^+ ∪ { 124 }
\end{verbatim}

(For small N the series becomes special due to the condition
den\textgreater0. For N=3 the series is den=1,2,3,4,5,124; den=1 means
Δτ=2π.) In the stage-1-only baseline, H = A*(z̄⊗z) (amplitude-weighted)
is used. This map is memoryless (one step is a pure function of z and
constants); the old hidden state wp and approximate σ estimation do not
exist.

\paragraph{2.3.1 Design Background (Hypothesis) of the Denominator
Series and
Its}\label{design-background-hypothesis-of-the-denominator-series-and-its}

Implementation Correspondence

The denominator series \{ N−2, N−1, N, N+1, N+2, 124 \} of Δτ = 2π/den
is not an invention of this chapter; it inherits, unchanged, the design
of the \textbf{denominator control experiments} of this series
(2026-09-03; design instruction document
\texttt{ChatGPT\_denominator\_controls\_N3\_N40\_mixedseed\_20260903/CLAUDE\_CODE\_RUN\_INSTRUCTION\_N3\_N40\_20260903.md}).
The background splits into two lines.

\textbf{(a) den=124 (fixed-clock control)} --- Inherited from the
step-size convention Δτ = 2π/L, \textbf{L=124} of the
interference-preserving dynamics series (e.g., line 27 \texttt{L=124} of
\texttt{.../program/pass1\_parents.py} in the qualification-and-seedless
series folder). It is the control with the clock fixed independently of
N (legacy convention).

\textbf{(b) den ∈ \{N−2..N+2\} (hypothesis series of a system-scaled
clock)} --- The hypothesis connects to the implementation through the
rotation angle of Eq. 20, ψ\_k = (2π/den)·σ\_k. Measuring σ\_max from
the parents' σ spectra saved in Chapter 1 (the \texttt{sigma} field of
the npz; aggregated in \texttt{analysis\_sweep\_summary\_v1.json},
\texttt{sigma\_max\_by\_N}):

{\def\LTcaptype{none} % do not increment counter
\begin{longtable}[]{@{}llll@{}}
\toprule\noalign{}
N & σ\_max & ψ\_max/2π (den=N) & ψ\_max/2π (den=124) \\
\midrule\noalign{}
\endhead
\bottomrule\noalign{}
\endlastfoot
3 & 1.414 & 0.471 & 0.011 \\
5 & 3.742 & 0.748 & 0.030 \\
10 & 8.928 & 0.893 & 0.072 \\
20 & 18.894 & 0.945 & 0.152 \\
30 & 28.898 & 0.963 & 0.233 \\
40 & 38.905 & 0.973 & 0.314 \\
\end{longtable}
}

σ\_max(N) grows nearly linearly with N (N − σ\_max ≈ 1.1). Therefore:

\begin{itemize}
\tightlist
\item
  In the \textbf{den ≈ N series}, the one-step rotation of the fastest
  eigenplane is placed at ψ\_max ≈ 2π·(1 − O(1)/N), i.e., in the
  \textbf{nearly-one-full-turn (near-commensurate with 2π, stroboscopic)
  regime}. The five points den = N, N±1, N±2 sweep the detuning from
  this commensurability.
\item
  With \textbf{den = 124}, all eigenplanes stay in the small-angle
  regime (ψ\_max/2π ≤ 0.314 within this sweep), a control close to the
  small-step approximation of the continuous flow dz/dt = Kz
  (\textbf{flow-like regime}).
\end{itemize}

\textbf{Implementation correspondence}: den is determined by the series
generation on line 38
(\texttt{pairs={[}(N+o,…)\ for\ o\ in\ OFFSETS\ if\ N+o\textgreater{}0{]}+{[}(124,\textquotesingle{}124\textquotesingle{}){]}})
and enters the dynamics only through \texttt{2.0*math.pi/den} on line
27. den is used nowhere else (as Eq. 20 states, the effect of den is
exhausted by the linear scaling of rotation angles and aliasing).

This chapter does not judge this design hypothesis. The den-dependent
observations are recorded as facts in §6 (denominator dependence of
onset; distribution of non-crossing runs).

\textbf{(Eq. 17) Stage 2: amplitude normalization} --- The input to the
generator construction is replaced by the unit-modulus z\^{}:

\begin{verbatim}
z^_e = e^{i θ_e} = e^{i arg z_e}
Ĥ_ef = A_ef · conj(z^_e) z^_f = A_ef · e^{i(θ_f − θ_e)}
\end{verbatim}

The generator thereby becomes a function of the \textbf{phases only}, as
in the old dynamics (Eq. 11). The amplitudes of the state z itself are
retained and continue to evolve as dynamical variables (the
normalization happens only inside the generator construction, at every
step).

\textbf{(Eq. 18) Real/imaginary decomposition of Ĥ} --- Ĥ is Hermitian
and decomposes uniquely into a real symmetric part and a real
antisymmetric part:

\begin{verbatim}
Ĥ = S + iK,
S_ef = A_ef cos(θ_f − θ_e)   (real symmetric)
K_ef = A_ef sin(θ_f − θ_e)   (real antisymmetric)
\end{verbatim}

This K is \textbf{the same matrix as the old generator (Eq. 11)}.

\textbf{(Eq. 19) Stage 3: extraction of the imaginary part (removal of
the cos symmetric part) and real orthogonal rotation} --- Instead of Ĥ,
the generator

\begin{verbatim}
H_3 = i K = i · Im(Ĥ)
\end{verbatim}

is used (i × real antisymmetric = Hermitian). Substituting into the
stage-1 map (Eq. 16):

\begin{verbatim}
z ← exp(−i Δτ · iK) z = exp(Δτ K) z
\end{verbatim}

exp(ΔτK) is the exponential of a real antisymmetric matrix, i.e., a
\textbf{real orthogonal matrix}. It therefore has the same conserved
quantities ‖z‖\^{}2, z\^{}Tz as the old dynamics (Eq. 15; same proof).
Since stage 2 makes K phase-only, the relative-equilibrium argument of
Eq. 15' also holds unchanged: along the parent orbit K is time-invariant
and the parent rotates rigidly.

\textbf{(Eq. 20) New eigenplane rotation angle} --- The eigenvalues w\_k
of H\_3 = iK are real; corresponding to ±iσ\_k of K, w = ±σ\_k. The
one-step rotation angle of eigenplane k is

\begin{verbatim}
ψ_k = Δτ · σ_k = (2π/den) · σ_k
\end{verbatim}

The rotation angle is \textbf{linear} in the eigenvalue, with no upper
bound. The phase factor e\^{}\{−iΔτw\} wraps around in w with period
2π/Δτ (\textbf{aliasing}): for eigenplanes with Δτσ\_k exceeding 2π, the
effective rotation angle is Δτσ\_k mod 2π.

\subsubsection{2.4 Mathematical Differences between the Old and New
Rotation Maps (the crux
of}\label{mathematical-differences-between-the-old-and-new-rotation-maps-the-crux-of}

this chapter's dynamics)

Both are ``real orthogonal rotations generated by the phase-only real
antisymmetric K'', but \textbf{the spectral map of rotation angles and
the clock differ mathematically}.

{\def\LTcaptype{none} % do not increment counter
\begin{longtable}[]{@{}
  >{\raggedright\arraybackslash}p{(\linewidth - 4\tabcolsep) * \real{0.3333}}
  >{\raggedright\arraybackslash}p{(\linewidth - 4\tabcolsep) * \real{0.3333}}
  >{\raggedright\arraybackslash}p{(\linewidth - 4\tabcolsep) * \real{0.3333}}@{}}
\toprule\noalign{}
\begin{minipage}[b]{\linewidth}\raggedright
Item
\end{minipage} & \begin{minipage}[b]{\linewidth}\raggedright
Old (Eqs. 12--14)
\end{minipage} & \begin{minipage}[b]{\linewidth}\raggedright
This chapter = stages 1+2+3 (Eqs. 16--20)
\end{minipage} \\
\midrule\noalign{}
\endhead
\bottomrule\noalign{}
\endlastfoot
Eigenplane rotation angle & φ\_k = 2 arctan(γ σ\_k/σ\^{}\_max) & ψ\_k =
(2π/den) σ\_k \\
Spectral dependence & nonlinear compression by arctan
(\textbar φ\textbar{} \textless{} π) & linear (with mod-2π aliasing) \\
Clock & σ\^{}\_max-normalized (fastest plane fixed at π/72;
system-intrinsic) & fixed Δτ = 2π/den (external; den is the sweep
parameter) \\
Obtaining σ & approximate σ\^{}\_max by 3 power iterations (warm start,
history-dependent) & exact eigenvalues via eigh (no estimation) \\
Realization of the map & Cayley rational form applied via Woodbury solve
(O(N\^{}3)/step) & exact exponential applied via spectral decomposition
(O(M\^{}3)/step) \\
Memory & carries hidden state wp (not a pure function) & memoryless
(pure function of z) \\
Conserved quantities & ‖z‖\^{}2, z\^{}Tz (real orthogonal) & same (real
orthogonal) \\
Relative equilibrium (parent) & holds (Eq. 15') & holds (consequence of
Eq. 19) \\
\end{longtable}
}

\textbf{(Eq. 21) Correspondence in the small-angle limit} --- In the
regime Δτσ\_k ≪ 1 and γσ\_k/σ\^{}\_max ≪ 1, arctan x ≈ x gives

\begin{verbatim}
φ_k ≈ (2γ/σ^_max) · σ_k = Δτ_eff · σ_k,   Δτ_eff = 2 tan(π/144)/σ^_max ≈ (π/72)/σ^_max
\end{verbatim}

That is, the old dynamics approximately coincides with ``Eq. 20 with
Δτ\_eff re-normalized by σ\^{}\_max at every step''. The differences
between the two reduce to (i) the value of Δτ (externally fixed 2π/den
vs.~the system-intrinsic (π/72)/σ\^{}\_max), (ii) the presence of arctan
compression and aliasing, and (iii) the exactness of σ. The measured
record of this correspondence is the last row of Appendix A (the
σ-normalized-clock reference experiment: late-segment slope 65.8
steps/decade vs.~the July canonical 64.0 steps/decade). The runs of this
chapter use the \textbf{fixed Δτ (stage 1)} throughout; the sweep of den
measures the effect of this clock difference.

\subsubsection{2.5 Measured Quantities and Conservation Laws ---
Definition of the Dormant
Fraction}\label{measured-quantities-and-conservation-laws-definition-of-the-dormant-fraction}

H⊥/H and Its Justification as the Inflation Indicator

\textbf{(Eq. 22) Readout plane} --- From the initial state z0 = v + δg
(after normalization):

\begin{verbatim}
p = Re z0 / ‖Re z0‖,   q' = Im z0 − (Im z0·p) p,   q = q'/‖q'‖
\end{verbatim}

p, q are \textbf{fixed} throughout the run (the Gram--Schmidt
orthonormal basis of the real 2-dimensional plane Π = span\_R(p, q)
spanned by the real and imaginary parts of the initial state). Since δ =
10\^{}-15, Π is essentially \textbf{the plane spanned by the parent v}.

\textbf{(Eq. 23) Measured quantities} --- For the state z at each step:

\begin{verbatim}
H⊥/H = ‖ z − p(p·z) − q(q·z) ‖^2 / ‖z‖^2   (energy fraction outside the readout plane Π)
H_total = z†z
closure = |z^Tz| / z†z
\end{verbatim}

\textbf{(Eq. 23') Why H⊥/H is the ``dormant fraction'' and does not pick
up the parent's motion} --- The orbit of the pure parent (δ=0) is a
relative equilibrium (Eqs. 15', 19): Z(t) = e\^{}\{−iσt\} Z(0).
Expanding,

\begin{verbatim}
e^{−iσt}(x + iy) = (x cos σt + y sin σt) + i (y cos σt − x sin σt),   x=Re Z(0), y=Im Z(0)
\end{verbatim}

i.e., the real and imaginary parts always remain within span\_R(x, y) =
Π. The complementary orthogonal projection z − p(p·z) − q(q·z) exactly
annihilates any complex combination inside Π, so \textbf{H⊥ is
identically 0 along the parent orbit}. Hence H⊥/H measures only
``everything \textbf{other than} the parent's rigid rotation'' ---
initially the out-of-Π component of the seed g (H⊥/H(0) ≈ δ\^{}2
\textasciitilde{} 10\^{}-30), and afterwards whatever the dynamics grows
from it. Because the seed is energy invisible at the parent's scale, the
July canonical calls this the \textbf{dormant fraction}.

\textbf{(Eq. 23'\,') Identity of definition with the July canonical} ---
The H⊥/H of this chapter is the same measurement as f(τ) of the July
canonical. Canonical program
\texttt{自発的分裂予備実験\_v1/run\_spontaneous\_splitting\_largeN\_v1.py}:

\begin{Shaded}
\begin{Highlighting}[]
    \DecValTok{57}\NormalTok{          Zp }\OperatorTok{=}\NormalTok{ Z }\OperatorTok{{-}}\NormalTok{ p }\OperatorTok{*}\NormalTok{ (p }\OperatorTok{@}\NormalTok{ Z) }\OperatorTok{{-}}\NormalTok{ q }\OperatorTok{*}\NormalTok{ (q }\OperatorTok{@}\NormalTok{ Z)}
    \DecValTok{58}\NormalTok{          htot }\OperatorTok{=} \BuiltInTok{float}\NormalTok{(np.real(np.conj(Z) }\OperatorTok{@}\NormalTok{ Z))}
    \DecValTok{59}\NormalTok{          f }\OperatorTok{=} \BuiltInTok{float}\NormalTok{(np.real(np.conj(Zp) }\OperatorTok{@}\NormalTok{ Zp)) }\OperatorTok{/}\NormalTok{ htot}
\end{Highlighting}
\end{Shaded}

(The construction of p, q is also identical; lines 46--48.) The curves
of this chapter therefore plot the same quantity with the same
definition as the July inflation figure
(dormant\_growth\_large\_n\_v1.png) and can be compared directly.

\textbf{Justification as the inflation indicator} --- By the above, the
time evolution of H⊥/H quantifies inflation-like development by three
points: ``(i) the initial value lies at the seed scale 10\^{}-30 (as
long as the parent is an equilibrium, there is no jump at the first
step); (ii) from there a straight line on the semilog plot --- i.e.,
constant-rate exponential amplification --- continues; (iii) it
saturates at O(10\^{}-2--1).'' When (i) fails (a jump to
10\textsuperscript{-8--10}-3 at the first step), that is a mismatch
injection and is distinguished from seed amplification (the deletion
controls of Appendix A are discriminated precisely by this (i)). The
onset (Eq. 24) is the indicator of reaching (iii).

Conservation: by Eq. 19 the map is real orthogonal, so both the
numerator \textbar z\^{}Tz\textbar{} and denominator z†z of closure are
theoretically invariant; the time variation of closure measures only the
accumulation of numerical rounding (measured in §6, G3). Norm
conservation guarantees H⊥/H ∈ {[}0,1{]}.

\textbf{(Eq. 24) onset (bookkeeping indicator)} --- The first step
number at which H⊥/H \textgreater{} 0.05 (−1 if none). A tally that does
not affect the dynamics.

\subsubsection{2.6 Correspondence with the Complex-Plane Readout ---
Design Intent of the
Inflation}\label{correspondence-with-the-complex-plane-readout-design-intent-of-the-inflation}

Figure and the Three Complex-Plane Figures

The complex-plane figures of Chapter 3 (step 0, final step, zoom into
the condensed center) are the readout paired with the inflation figure
(H⊥/H curves) of this chapter. The identity that fixes their relation is
given here first.

\textbf{Caution (two easily confused ``planes'')} --- The complex plane
of the figures is ``the plane on which the M per-edge values z\_e ∈ C
are plotted'', whereas the parent plane Π = span\_R(p,q), the reference
of H⊥/H, is ``a real 2-dimensional subspace of the state space C\^{}M''.
They are distinct, but connected by the following identity.

\textbf{(Eq. 25) Per-edge representation of states inside Π} --- With x
= Re z0, y = Im z0, v ≈ z0 (δ=10\^{}-15), x = (v+v̄)/2 and y = (v−v̄)/(2i)
give

\begin{verbatim}
z ∈ span_C{x, y} = Π  <=>  z_e = a·v_e + b·conj(v_e)   (a, b ∈ C constants independent of the edge)
\end{verbatim}

That is, **motion confined to Π can, in the figure, produce only
deformations within the two-complex-parameter family ``uniform
rotation/scaling of the parent's star (a·v\_e) plus admixture of its
mirror image (b·v̄\_e)``\textbf{. In particular the pure parent orbit
Z(t)=e\^{}\{−iσt\}Z(0) (Eq. 15') has a=e\^{}\{−iσt\}, b=0: }the whole
star rotates uniformly about the origin without changing its shape**.
Conversely, the appearance in the figure of shapes not expressible in
this family (e.g., phase dispersal in all directions = a ring) is the
configurational expression of the state having left Π.

By this identity, the three complex-plane figures share the roles of
corroborating, from the configuration side, the start point, end point,
and interior of the end point of the inflation figure (H⊥/H curves):

{\def\LTcaptype{none} % do not increment counter
\begin{longtable}[]{@{}
  >{\raggedright\arraybackslash}p{(\linewidth - 4\tabcolsep) * \real{0.3333}}
  >{\raggedright\arraybackslash}p{(\linewidth - 4\tabcolsep) * \real{0.3333}}
  >{\raggedright\arraybackslash}p{(\linewidth - 4\tabcolsep) * \real{0.3333}}@{}}
\toprule\noalign{}
\begin{minipage}[b]{\linewidth}\raggedright
Figure
\end{minipage} & \begin{minipage}[b]{\linewidth}\raggedright
Correspondence to the H⊥/H curve
\end{minipage} & \begin{minipage}[b]{\linewidth}\raggedright
What it reads
\end{minipage} \\
\midrule\noalign{}
\endhead
\bottomrule\noalign{}
\endlastfoot
step-0 figure & start of the curve (H⊥/H ≈ 10\^{}-31) & the contents of
Π = the parent's shape (two antipodal pairs, 4-bundle star). That the
step-1 figure is indistinguishable from step 0 visualizes f(1) staying
at seed scale (§2.5 (i)) \\
final-step figure & saturation region of the curve (H⊥/H
\textasciitilde{} 0.05--0.46) & departure to shapes not expressible in
the family of Eq. 25 (star → ring) = configuration in which out-of-Π
components dominate. onset (Eq. 24) roughly corresponds to when the
shape breakdown becomes visible \\
zoom into the condensed center & complement of what the curve does not
say & H⊥/H measures only ``how much left Π''. The zoom resolves ``the
organization at the destination'' --- whether angular clusters are
condensation onto exactly identical complex values (exact degeneracy) or
bundles of finite width \\
\end{longtable}
}

The implementation of the figures (grid rendering, cluster-extraction
algorithm, line-number correspondence) is described in Chapter 3.

\subsection{3. Implementation Method}\label{implementation-method}

\begin{itemize}
\tightlist
\item
  The sweep body \texttt{run\_N3\_N40\_stage123\_v1.py} is
  self-contained (no engine import). The bundled engine
  \texttt{run\_n\_scaling\_lowrank\_v1.py} serves only as the definition
  body of §2.2 (the old mathematics) and for the parent generation of
  Chapter 1; it plays no role in this chapter's dynamics loop.
\item
  This program is derived from the ChatGPT-authored canonical
  \texttt{run\_and\_plot\_N3\_N40\_mixedseed\_20260903.py} through
  \textbf{a documented chain of minimal diffs only} (Appendix B); the
  dynamics function \texttt{one\_step} is identical to the stage-3
  version that passed the three gates in the N=40 single-factor
  experiments.
\item
  The initial data are the \texttt{Z0} of the Chapter-1 static-parent
  npz files, loaded per N (no regeneration, no modification).
\item
  The input gate \texttt{check\_sweep\_inputs\_v1.py} verifies post hoc
  that the stored \texttt{Z{[}0{]}} of all 228 runs is bit-identical to
  the static parents' \texttt{Z0}.
\item
  The numbers in §6 are transcribed from the output of the aggregation
  program \texttt{analyze\_sweep\_summary\_v1.py},
  \texttt{results/analysis\_sweep\_summary\_v1.json} (no manual
  computation or manual tallying).
\end{itemize}

\textbf{Stage ⇔ equation ⇔ implementation correspondence table}
(implementation: \texttt{run\_N3\_N40\_stage123\_v1.py}):

{\def\LTcaptype{none} % do not increment counter
\begin{longtable}[]{@{}
  >{\raggedright\arraybackslash}p{(\linewidth - 6\tabcolsep) * \real{0.2500}}
  >{\raggedright\arraybackslash}p{(\linewidth - 6\tabcolsep) * \real{0.2500}}
  >{\raggedright\arraybackslash}p{(\linewidth - 6\tabcolsep) * \real{0.2500}}
  >{\raggedright\arraybackslash}p{(\linewidth - 6\tabcolsep) * \real{0.2500}}@{}}
\toprule\noalign{}
\begin{minipage}[b]{\linewidth}\raggedright
Stage
\end{minipage} & \begin{minipage}[b]{\linewidth}\raggedright
Content
\end{minipage} & \begin{minipage}[b]{\linewidth}\raggedright
Equations
\end{minipage} & \begin{minipage}[b]{\linewidth}\raggedright
Implementation
\end{minipage} \\
\midrule\noalign{}
\endhead
\bottomrule\noalign{}
\endlastfoot
Stage 1 & explicit matrix + exact spectral map & Eq. 16 & lines 21--22
(\texttt{H\_of}), lines 27--28 (\texttt{eigh} and application of the
phase factor) \\
Stage 1 & fixed Δτ=2π/den clock and denominator series & Eq. 16 &
\texttt{2.0*math.pi/den} on line 27; line 38 (series generation) \\
Stage 2 & amplitude normalization (z\textsuperscript{=e}\{iθ\}) & Eq. 17
& \texttt{np.exp(1j*np.angle(z))} on line 26 \\
Stage 3 & imaginary-part extraction H\_3=i·Im(Ĥ) → real orthogonal
rotation & Eqs. 18, 19 & \texttt{(1j*np.imag(H))} on line 26 \\
--- & readout plane, measured quantities, onset & Eqs. 22, 23, 24 &
lines 29--30, 31--32, 45--46 \\
\end{longtable}
}

(Contrast: the old Eqs. 12--14 exist only in engine lines 122--132 and
134--141 and are never called in this chapter's loop.)

\subsection{4. Detailed Design}\label{detailed-design}

\subsubsection{4.1 Overall Flow}\label{overall-flow}

\begin{verbatim}
for N in 3..40:
  [initialization]  load static parent Z0 → adjacency A (Eq. 10) → readout plane p,q (Eq. 22) → denominator series (Eq. 16)
  for den in series:
    [loop]  t=0..500: record state and measured quantities (Eq. 23); apply one_step (Eqs. 16–19) if t<500
    [finalization (den)]  save state npz; append onset (Eq. 24) etc. to the summary row
[finalization (all)]  write timeseries/summary CSVs → 8×5 grid figure → RUN_METADATA
\end{verbatim}

\subsubsection{4.2 Overall Data Flow}\label{overall-data-flow}

\begin{itemize}
\tightlist
\item
  \textbf{Input}: \texttt{Z0} from
  \texttt{parents/parent\_static\_N\{N:05d\}\_makeparent\_20260905.npz}
  (38 files; Chapter-1 products; SHA256 canonical in SHA256SUMS.txt)
\item
  \textbf{Parameters}:

  \begin{itemize}
  \tightlist
  \item
    \texttt{STEPS\ =\ 500} \ldots{} steps per run (no early stopping;
    fixed)
  \item
    \texttt{OFFSETS\ =\ (-2,-1,0,1,2)} \ldots{} offsets of the
    denominator series from N (Eq. 16)
  \item
    denominator 124 \ldots{} the control denominator common to all N
  \item
    dtype: state complex128 / real float64 (fixed by the assert on line
    12)
  \end{itemize}
\item
  \textbf{Output}:

  \begin{itemize}
  \tightlist
  \item
    \texttt{results/hm\_N\{N\}\_den\_\{den\}\_states\_500.npz} × 228
    (\texttt{Z}(501×M), \texttt{N}, \texttt{denominator},
    \texttt{steps})
  \item
    \texttt{results/timeseries\_64bit\_with124\_N3\_N40.csv} (columns:
    N, series, denominator, step, Hperp\_frac, H\_total,
    global\_closure; 228×501 = 114,228 rows)
  \item
    \texttt{results/summary\_64bit\_with124\_N3\_N40.csv} (columns: N,
    series, denominator, onset\_gt\_0.05, initial, step1, final, max;
    228 rows)
  \item
    \texttt{results/fig\_Hperp\_denominator\_controls\_with\_124\_N3\_N40\_stage123.png}
    (target figure)
  \item
    \texttt{results/RUN\_METADATA\_N3\_N40\_stage123.json}
  \end{itemize}
\end{itemize}

\subsubsection{4.3 Individual Processes}\label{individual-processes}

\paragraph{\texorpdfstring{4.3.1 Initialization
(\texttt{run\_N3\_N40\_stage123\_v1.py})}{4.3.1 Initialization (run\_N3\_N40\_stage123\_v1.py)}}\label{initialization-run_n3_n40_stage123_v1.py}

Constants and dtype fixing:

\begin{Shaded}
\begin{Highlighting}[]
    \DecValTok{11}\NormalTok{  STEPS}\OperatorTok{=}\DecValTok{500}\OperatorTok{;}\NormalTok{ OFFSETS}\OperatorTok{=}\NormalTok{(}\OperatorTok{{-}}\DecValTok{2}\NormalTok{,}\OperatorTok{{-}}\DecValTok{1}\NormalTok{,}\DecValTok{0}\NormalTok{,}\DecValTok{1}\NormalTok{,}\DecValTok{2}\NormalTok{)}
    \DecValTok{12}  \ControlFlowTok{assert}\NormalTok{ np.dtype(np.float64).itemsize}\OperatorTok{==}\DecValTok{8} \KeywordTok{and}\NormalTok{ np.dtype(np.complex128).itemsize}\OperatorTok{==}\DecValTok{16}
\end{Highlighting}
\end{Shaded}

Edges and adjacency (Eqs. 1, 10):

\begin{Shaded}
\begin{Highlighting}[]
    \DecValTok{14}  \KeywordTok{def}\NormalTok{ edges(N):}
    \DecValTok{15}\NormalTok{      a,b}\OperatorTok{=}\NormalTok{np.triu\_indices(N,k}\OperatorTok{=}\DecValTok{1}\NormalTok{)}\OperatorTok{;} \ControlFlowTok{return}\NormalTok{ a.astype(np.int64),b.astype(np.int64)}
    \DecValTok{16}  \KeywordTok{def}\NormalTok{ adjacency(N):}
    \DecValTok{17}\NormalTok{      ea,eb}\OperatorTok{=}\NormalTok{edges(N)}\OperatorTok{;}\NormalTok{ M}\OperatorTok{=}\BuiltInTok{len}\NormalTok{(ea)}\OperatorTok{;}\NormalTok{ A}\OperatorTok{=}\NormalTok{np.zeros((M,M),dtype}\OperatorTok{=}\NormalTok{np.float64)}
    \DecValTok{18}      \ControlFlowTok{for}\NormalTok{ e }\KeywordTok{in} \BuiltInTok{range}\NormalTok{(M):}
    \DecValTok{19}\NormalTok{          share}\OperatorTok{=}\NormalTok{(ea}\OperatorTok{==}\NormalTok{ea[e])}\OperatorTok{|}\NormalTok{(ea}\OperatorTok{==}\NormalTok{eb[e])}\OperatorTok{|}\NormalTok{(eb}\OperatorTok{==}\NormalTok{ea[e])}\OperatorTok{|}\NormalTok{(eb}\OperatorTok{==}\NormalTok{eb[e])}\OperatorTok{;}\NormalTok{ share[e]}\OperatorTok{=}\VariableTok{False}\OperatorTok{;}\NormalTok{ A[e,share]}\OperatorTok{=}\FloatTok{1.0}
    \DecValTok{20}      \ControlFlowTok{return}\NormalTok{ A}
\end{Highlighting}
\end{Shaded}

Parent loading, readout plane (Eq. 22), denominator series (Eq. 16):

\begin{Shaded}
\begin{Highlighting}[]
    \DecValTok{34}  \ControlFlowTok{for}\NormalTok{ N }\KeywordTok{in} \BuiltInTok{range}\NormalTok{(}\DecValTok{3}\NormalTok{,}\DecValTok{41}\NormalTok{):}
    \DecValTok{35}      \CommentTok{\# 初期データ: 各 N の静的親ファイルの Z0 を使用}
    \DecValTok{36}\NormalTok{      z0}\OperatorTok{=}\NormalTok{np.array(np.load(os.path.join(PARENT\_DIR,}\SpecialStringTok{f\textquotesingle{}parent\_static\_N}\SpecialCharTok{\{}\NormalTok{N}\SpecialCharTok{:05d\}}\SpecialStringTok{\_makeparent\_20260905.npz\textquotesingle{}}\NormalTok{))[}\StringTok{\textquotesingle{}Z0\textquotesingle{}}\NormalTok{],dtype}\OperatorTok{=}\NormalTok{np.complex128,copy}\OperatorTok{=}\VariableTok{True}\NormalTok{)}
    \DecValTok{37}\NormalTok{      A}\OperatorTok{=}\NormalTok{adjacency(N)}\OperatorTok{;}\NormalTok{ p,q}\OperatorTok{=}\NormalTok{plane(z0)}
    \DecValTok{38}\NormalTok{      pairs}\OperatorTok{=}\NormalTok{[(N}\OperatorTok{+}\NormalTok{o, }\SpecialStringTok{f\textquotesingle{}N}\SpecialCharTok{\{}\NormalTok{o}\SpecialCharTok{:+d\}}\SpecialStringTok{\textquotesingle{}} \ControlFlowTok{if}\NormalTok{ o }\ControlFlowTok{else} \StringTok{\textquotesingle{}N\textquotesingle{}}\NormalTok{) }\ControlFlowTok{for}\NormalTok{ o }\KeywordTok{in}\NormalTok{ OFFSETS }\ControlFlowTok{if}\NormalTok{ N}\OperatorTok{+}\NormalTok{o}\OperatorTok{\textgreater{}}\DecValTok{0}\NormalTok{] }\OperatorTok{+}\NormalTok{ [(}\DecValTok{124}\NormalTok{,}\StringTok{\textquotesingle{}124\textquotesingle{}}\NormalTok{)]}
\end{Highlighting}
\end{Shaded}

\begin{Shaded}
\begin{Highlighting}[]
    \DecValTok{29}  \KeywordTok{def}\NormalTok{ plane(v):}
    \DecValTok{30}\NormalTok{      p}\OperatorTok{=}\NormalTok{v.real.astype(np.float64,copy}\OperatorTok{=}\VariableTok{True}\NormalTok{)}\OperatorTok{;}\NormalTok{ p}\OperatorTok{/=}\NormalTok{np.linalg.norm(p)}\OperatorTok{;}\NormalTok{ q}\OperatorTok{=}\NormalTok{v.imag.astype(np.float64,copy}\OperatorTok{=}\VariableTok{True}\NormalTok{)}\OperatorTok{;}\NormalTok{ q}\OperatorTok{{-}=}\NormalTok{np.dot(q,p)}\OperatorTok{*}\NormalTok{p}\OperatorTok{;}\NormalTok{ q}\OperatorTok{/=}\NormalTok{np.linalg.norm(q)}\OperatorTok{;} \ControlFlowTok{return}\NormalTok{ p,q}
\end{Highlighting}
\end{Shaded}

\paragraph{4.3.2 Loop (the one-step dynamics --- Eqs.
16--20)}\label{loop-the-one-step-dynamics-eqs.-1620}

\begin{Shaded}
\begin{Highlighting}[]
    \DecValTok{21}  \KeywordTok{def}\NormalTok{ H\_of(z,A):}
    \DecValTok{22}\NormalTok{      H}\OperatorTok{=}\NormalTok{A}\OperatorTok{*}\NormalTok{(np.conj(z)[:,}\VariableTok{None}\NormalTok{]}\OperatorTok{*}\NormalTok{z[}\VariableTok{None}\NormalTok{,:])}\OperatorTok{;}\NormalTok{ np.fill\_diagonal(H,}\FloatTok{0.0}\NormalTok{)}\OperatorTok{;} \ControlFlowTok{return}\NormalTok{ H.astype(np.complex128,copy}\OperatorTok{=}\VariableTok{False}\NormalTok{)}
    \DecValTok{23}  \KeywordTok{def}\NormalTok{ one\_step(z,A,den):}
    \DecValTok{24}      \CommentTok{\# 段3の最小変更（唯一の力学変更点）: 位相のみ生成子 Ĥ の虚部だけを取る H=i·K（K=sin(Δθ) 実反対称）。}
    \DecValTok{25}      \CommentTok{\# exp({-}iΔτ·iK)=exp(Δτ·K) の実直交回転となり、Z\^{}T Z（零閉塞）と ‖Z‖ を厳密保存する。}
    \DecValTok{26}\NormalTok{      H}\OperatorTok{=}\NormalTok{H\_of(np.exp(}\OtherTok{1j}\OperatorTok{*}\NormalTok{np.angle(z)),A)}\OperatorTok{;}\NormalTok{ H}\OperatorTok{=}\NormalTok{(}\OtherTok{1j}\OperatorTok{*}\NormalTok{np.imag(H)).astype(np.complex128,copy}\OperatorTok{=}\VariableTok{False}\NormalTok{)}
    \DecValTok{27}\NormalTok{      w,V}\OperatorTok{=}\NormalTok{np.linalg.eigh(H)}\OperatorTok{;}\NormalTok{ phase}\OperatorTok{=}\NormalTok{np.exp(}\OperatorTok{{-}}\OtherTok{1j}\OperatorTok{*}\NormalTok{np.float64(}\FloatTok{2.0}\OperatorTok{*}\NormalTok{math.pi}\OperatorTok{/}\NormalTok{den)}\OperatorTok{*}\NormalTok{w)}
    \DecValTok{28}      \ControlFlowTok{return}\NormalTok{ (V}\OperatorTok{@}\NormalTok{(phase}\OperatorTok{*}\NormalTok{(V.conj().T}\OperatorTok{@}\NormalTok{z))).astype(np.complex128,copy}\OperatorTok{=}\VariableTok{False}\NormalTok{)}
\end{Highlighting}
\end{Shaded}

\begin{itemize}
\tightlist
\item
  First half of line 26, \texttt{np.exp(1j*np.angle(z))}: \textbf{stage
  2} = Eq. 17 (amplitude normalization; the generator input is replaced
  by z\textsuperscript{=e}\{iθ\}).
\item
  Second half of line 26, \texttt{(1j*np.imag(H))}: \textbf{stage 3} =
  the imaginary-part extraction of Eq. 18 → the generator H\_3=i·K of
  Eq. 19 (\texttt{np.imag} is the elementwise imaginary part; the
  imaginary part of the Hermitian Ĥ is the real antisymmetric K).
\item
  Lines 27--28: \textbf{stage 1} = Eq. 16 (exact spectral decomposition
  by \texttt{np.linalg.eigh}; application of the phase factor
  e\^{}\{−i(2π/den)w\} with fixed Δτ=2π/den). The eigenplane rotation
  angle is Eq. 20, ψ\_k=(2π/den)σ\_k, which differs from the old Eq. 14
  (arctan compression, σ\^{}\_max normalization) (§2.4).
\end{itemize}

Measurement and recording (Eq. 23):

\begin{Shaded}
\begin{Highlighting}[]
    \DecValTok{31}  \KeywordTok{def}\NormalTok{ metrics(z,p,q):}
    \DecValTok{32}\NormalTok{      h}\OperatorTok{=}\NormalTok{np.vdot(z,z).real}\OperatorTok{;}\NormalTok{ zp}\OperatorTok{=}\NormalTok{z}\OperatorTok{{-}}\NormalTok{p}\OperatorTok{*}\NormalTok{np.dot(p,z)}\OperatorTok{{-}}\NormalTok{q}\OperatorTok{*}\NormalTok{np.dot(q,z)}\OperatorTok{;}\NormalTok{ hp}\OperatorTok{=}\NormalTok{np.vdot(zp,zp).real}\OperatorTok{;} \ControlFlowTok{return} \BuiltInTok{float}\NormalTok{(hp}\OperatorTok{/}\NormalTok{h),}\BuiltInTok{float}\NormalTok{(h),}\BuiltInTok{float}\NormalTok{(}\BuiltInTok{abs}\NormalTok{(z}\OperatorTok{@}\NormalTok{z)}\OperatorTok{/}\NormalTok{h)}
\end{Highlighting}
\end{Shaded}

\begin{Shaded}
\begin{Highlighting}[]
    \DecValTok{40}\NormalTok{          z}\OperatorTok{=}\NormalTok{z0.copy()}\OperatorTok{;}\NormalTok{ vals}\OperatorTok{=}\NormalTok{np.empty(STEPS}\OperatorTok{+}\DecValTok{1}\NormalTok{,np.float64)}\OperatorTok{;}\NormalTok{ states}\OperatorTok{=}\NormalTok{np.empty((STEPS}\OperatorTok{+}\DecValTok{1}\NormalTok{,z.size),np.complex128)}\OperatorTok{;}\NormalTok{ closures}\OperatorTok{=}\NormalTok{np.empty(STEPS}\OperatorTok{+}\DecValTok{1}\NormalTok{,np.float64)}\OperatorTok{;}\NormalTok{ htot}\OperatorTok{=}\NormalTok{np.empty(STEPS}\OperatorTok{+}\DecValTok{1}\NormalTok{,np.float64)}
    \DecValTok{41}          \ControlFlowTok{for}\NormalTok{ t }\KeywordTok{in} \BuiltInTok{range}\NormalTok{(STEPS}\OperatorTok{+}\DecValTok{1}\NormalTok{):}
    \DecValTok{42}\NormalTok{              states[t]}\OperatorTok{=}\NormalTok{z}\OperatorTok{;}\NormalTok{ vals[t],htot[t],closures[t]}\OperatorTok{=}\NormalTok{metrics(z,p,q)}
    \DecValTok{43}              \ControlFlowTok{if}\NormalTok{ t}\OperatorTok{\textless{}}\NormalTok{STEPS: z}\OperatorTok{=}\NormalTok{one\_step(z,A,den)}
\end{Highlighting}
\end{Shaded}

\textbf{Stopping conditions and exceptional behavior (elements not
formulated as equations)}: - There is no early stopping. Every run
consists of a \textbf{fixed 501} recordings (t=0..500) and 500
applications of the map (lines 41--43). onset (Eq. 24) is used only for
tallying and never stops the dynamics (the
\texttt{np.flatnonzero(vals\textgreater{}0.05)} on lines 45--46). - For
runs at larger N, numpy RuntimeWarnings (divide by zero / overflow /
invalid value encountered in matmul) may be printed for the matrix
product \texttt{V@(phase*(V†z))}. These are warnings, not exceptions;
execution does not stop. In this series, bit-identity of reruns on
identical inputs (Chapter 1 G1 and this chapter's input gate) confirms
they do not affect the results. - For N=3,4 the
\texttt{N+o\textgreater{}0} filter of Eq. 16 admits small denominators
such as den=1 (Δτ=2π) into the series. There is no special-casing in the
program (a consequence of the comprehension on line 38). From the
viewpoint of Eq. 20, small den means large Δτ --- the regime where
aliasing (Δτσ\_k mod 2π) acts strongly.

\paragraph{4.3.3 Finalization (saving, plotting,
metadata)}\label{finalization-saving-plotting-metadata}

\begin{Shaded}
\begin{Highlighting}[]
    \DecValTok{44}\NormalTok{          np.savez\_compressed(os.path.join(OUT,}\SpecialStringTok{f\textquotesingle{}hm\_N}\SpecialCharTok{\{}\NormalTok{N}\SpecialCharTok{\}}\SpecialStringTok{\_den\_}\SpecialCharTok{\{}\NormalTok{den}\SpecialCharTok{\}}\SpecialStringTok{\_states\_500.npz\textquotesingle{}}\NormalTok{),Z}\OperatorTok{=}\NormalTok{states,N}\OperatorTok{=}\NormalTok{np.int64(N),denominator}\OperatorTok{=}\NormalTok{np.int64(den),steps}\OperatorTok{=}\NormalTok{np.int64(STEPS))}
    \DecValTok{45}\NormalTok{          rows.extend((N,label,den,t,vals[t],htot[t],closures[t]) }\ControlFlowTok{for}\NormalTok{ t }\KeywordTok{in} \BuiltInTok{range}\NormalTok{(STEPS}\OperatorTok{+}\DecValTok{1}\NormalTok{))}\OperatorTok{;}\NormalTok{ ix}\OperatorTok{=}\NormalTok{np.flatnonzero(vals}\OperatorTok{\textgreater{}}\FloatTok{0.05}\NormalTok{)}
    \DecValTok{46}\NormalTok{          summaries.append((N,label,den,}\BuiltInTok{int}\NormalTok{(ix[}\DecValTok{0}\NormalTok{]) }\ControlFlowTok{if}\NormalTok{ ix.size }\ControlFlowTok{else} \OperatorTok{{-}}\DecValTok{1}\NormalTok{,}\BuiltInTok{float}\NormalTok{(vals[}\DecValTok{0}\NormalTok{]),}\BuiltInTok{float}\NormalTok{(vals[}\DecValTok{1}\NormalTok{]),}\BuiltInTok{float}\NormalTok{(vals[}\OperatorTok{{-}}\DecValTok{1}\NormalTok{]),}\BuiltInTok{float}\NormalTok{(vals.}\BuiltInTok{max}\NormalTok{())))}
\end{Highlighting}
\end{Shaded}

The CSVs, figure, and metadata (lines 48--66) only write out measured
values and render the figure; they do not affect the dynamics. The
figure is an 8×5 grid (lines 55--64), 38 panels used and 2 panels off
(line 64).

\paragraph{\texorpdfstring{4.3.4 Input Gate
(\texttt{check\_sweep\_inputs\_v1.py})}{4.3.4 Input Gate (check\_sweep\_inputs\_v1.py)}}\label{input-gate-check_sweep_inputs_v1.py}

\begin{Shaded}
\begin{Highlighting}[]
    \DecValTok{18}  \ControlFlowTok{for}\NormalTok{ N }\KeywordTok{in} \BuiltInTok{range}\NormalTok{(}\DecValTok{3}\NormalTok{, }\DecValTok{41}\NormalTok{):}
    \DecValTok{19}\NormalTok{      Z0 }\OperatorTok{=}\NormalTok{ np.load(os.path.join(PARENT\_DIR, }\SpecialStringTok{f\textquotesingle{}parent\_static\_N}\SpecialCharTok{\{}\NormalTok{N}\SpecialCharTok{:05d\}}\SpecialStringTok{\_makeparent\_20260905.npz\textquotesingle{}}\NormalTok{))[}\StringTok{\textquotesingle{}Z0\textquotesingle{}}\NormalTok{]}
    \DecValTok{20}\NormalTok{      dens }\OperatorTok{=}\NormalTok{ [N }\OperatorTok{+}\NormalTok{ o }\ControlFlowTok{for}\NormalTok{ o }\KeywordTok{in}\NormalTok{ (}\OperatorTok{{-}}\DecValTok{2}\NormalTok{, }\OperatorTok{{-}}\DecValTok{1}\NormalTok{, }\DecValTok{0}\NormalTok{, }\DecValTok{1}\NormalTok{, }\DecValTok{2}\NormalTok{) }\ControlFlowTok{if}\NormalTok{ N }\OperatorTok{+}\NormalTok{ o }\OperatorTok{\textgreater{}} \DecValTok{0}\NormalTok{] }\OperatorTok{+}\NormalTok{ [}\DecValTok{124}\NormalTok{]}
    \DecValTok{21}      \ControlFlowTok{for}\NormalTok{ den }\KeywordTok{in}\NormalTok{ dens:}
    \DecValTok{22}\NormalTok{          p }\OperatorTok{=}\NormalTok{ os.path.join(RESULT\_DIR, }\SpecialStringTok{f\textquotesingle{}hm\_N}\SpecialCharTok{\{}\NormalTok{N}\SpecialCharTok{\}}\SpecialStringTok{\_den\_}\SpecialCharTok{\{}\NormalTok{den}\SpecialCharTok{\}}\SpecialStringTok{\_states\_500.npz\textquotesingle{}}\NormalTok{)}
    \DecValTok{23}\NormalTok{          same }\OperatorTok{=} \BuiltInTok{bool}\NormalTok{(np.array\_equal(np.load(p)[}\StringTok{\textquotesingle{}Z\textquotesingle{}}\NormalTok{][}\DecValTok{0}\NormalTok{], Z0))}
    \DecValTok{24}\NormalTok{          n\_checked }\OperatorTok{+=} \DecValTok{1}
    \DecValTok{25}          \ControlFlowTok{if} \KeywordTok{not}\NormalTok{ same:}
    \DecValTok{26}              \BuiltInTok{print}\NormalTok{(}\SpecialStringTok{f\textquotesingle{}MISMATCH: N=}\SpecialCharTok{\{}\NormalTok{N}\SpecialCharTok{\}}\SpecialStringTok{ den=}\SpecialCharTok{\{}\NormalTok{den}\SpecialCharTok{\}}\SpecialStringTok{\textquotesingle{}}\NormalTok{)}
    \DecValTok{27}\NormalTok{              ok }\OperatorTok{=} \VariableTok{False}
\NormalTok{...}
    \DecValTok{35}\NormalTok{  sys.exit(}\DecValTok{0} \ControlFlowTok{if}\NormalTok{ ok }\ControlFlowTok{else} \DecValTok{1}\NormalTok{)}
\end{Highlighting}
\end{Shaded}

It verifies that the \texttt{Z{[}0{]}} of every stored npz is
\texttt{np.array\_equal} (bit-identical) to the corresponding static
parent \texttt{Z0}; a single mismatch yields exit code 1.

\subsection{5. Execution Results}\label{execution-results}

\subsubsection{5.1 Reproduction Commands}\label{reproduction-commands}

\begin{Shaded}
\begin{Highlighting}[]
\BuiltInTok{cd}\NormalTok{ N3\_N40\_stage123\_sweep\_20260905}
\ExtensionTok{python3}\NormalTok{ run\_N3\_N40\_stage123\_v1.py        }\CommentTok{\# sweep body}
\ExtensionTok{python3}\NormalTok{ check\_sweep\_inputs\_v1.py         }\CommentTok{\# input gate}
\ExtensionTok{python3}\NormalTok{ analyze\_sweep\_summary\_v1.py      }\CommentTok{\# aggregation for the numbers in §6}
\CommentTok{\# or ./run\_all.sh (parents → sweep → gate → aggregation → plots)}
\end{Highlighting}
\end{Shaded}

\subsubsection{5.2 Execution Environment}\label{execution-environment}

\begin{itemize}
\tightlist
\item
  Python 3.9.6 (\texttt{.venv/bin/python3}), numpy 2.0.2 (BLAS/LAPACK:
  macOS Accelerate), matplotlib (grid rendering)
\item
  macOS 26.3.1 (arm64)
\end{itemize}

\subsubsection{5.3 Execution Time}\label{execution-time}

About \textbf{45--50 minutes} for the sweep body (38 N × 6 denominators
× 500 steps, including full state saving; measured 2026-09-05; dominated
by the M×M \texttt{eigh} × 500 × 6; N=40 alone about 3 minutes). Input
gate and aggregation: under 1 minute each.

\subsubsection{5.4 Verification Gates}\label{verification-gates}

{\def\LTcaptype{none} % do not increment counter
\begin{longtable}[]{@{}
  >{\raggedright\arraybackslash}p{(\linewidth - 6\tabcolsep) * \real{0.2500}}
  >{\raggedright\arraybackslash}p{(\linewidth - 6\tabcolsep) * \real{0.2500}}
  >{\raggedright\arraybackslash}p{(\linewidth - 6\tabcolsep) * \real{0.2500}}
  >{\raggedright\arraybackslash}p{(\linewidth - 6\tabcolsep) * \real{0.2500}}@{}}
\toprule\noalign{}
\begin{minipage}[b]{\linewidth}\raggedright
Gate
\end{minipage} & \begin{minipage}[b]{\linewidth}\raggedright
Pass condition
\end{minipage} & \begin{minipage}[b]{\linewidth}\raggedright
Measured
\end{minipage} & \begin{minipage}[b]{\linewidth}\raggedright
Verdict
\end{minipage} \\
\midrule\noalign{}
\endhead
\bottomrule\noalign{}
\endlastfoot
G1: input identity & \texttt{Z{[}0{]}} of all 228 npz bit-identical to
the corresponding static parent \texttt{Z0} & checked 228 / MISMATCH 0 &
\textbf{PASS} \\
G2: completion & \texttt{ALL\ DONE} printed, exit code 0 & confirmed &
\textbf{PASS} \\
G3: conserved quantity (observation) & closure =
\textbar z\^{}Tz\textbar/H stays small over all runs and steps & step-0
max 5.19e-15; global max 3.72e-13 & \textbf{PASS} \\
\end{longtable}
}

(The G3 numbers are \texttt{global\_closure\_step0\_max} /
\texttt{global\_closure\_all\_max} of
\texttt{analysis\_sweep\_summary\_v1.json}.)

\subsubsection{5.5 Data}\label{data}

{\def\LTcaptype{none} % do not increment counter
\begin{longtable}[]{@{}
  >{\raggedright\arraybackslash}p{(\linewidth - 2\tabcolsep) * \real{0.5000}}
  >{\raggedright\arraybackslash}p{(\linewidth - 2\tabcolsep) * \real{0.5000}}@{}}
\toprule\noalign{}
\begin{minipage}[b]{\linewidth}\raggedright
Item
\end{minipage} & \begin{minipage}[b]{\linewidth}\raggedright
Content
\end{minipage} \\
\midrule\noalign{}
\endhead
\bottomrule\noalign{}
\endlastfoot
Folder & \texttt{N3\_N40\_stage123\_sweep\_20260905/results/} \\
State npz & \texttt{hm\_N\{N\}\_den\_\{den\}\_states\_500.npz} × 228
(tens of KB to \textasciitilde5.9MB each; actual sizes per
SHA256SUMS.txt and the bundle) \\
Timeseries CSV & \texttt{timeseries\_64bit\_with124\_N3\_N40.csv}
(\textasciitilde8.7MB, 114,228 rows) \\
Summary CSV & \texttt{summary\_64bit\_with124\_N3\_N40.csv} (228
rows) \\
Metadata & \texttt{RUN\_METADATA\_N3\_N40\_stage123.json} \\
Aggregation & \texttt{analysis\_sweep\_summary\_v1.json} (source of the
numbers in §6) \\
SHA256 & Canonical values for all files in the bundled
\texttt{SHA256SUMS.txt}. Sweep program:
\texttt{1abf2353fee2e4f56f05e7a6f149fd086885136beb61ab571b48a56b09691567\ \ run\_N3\_N40\_stage123\_v1.py} \\
\end{longtable}
}

\subsubsection{5.6 Figures}\label{figures}

\begin{itemize}
\tightlist
\item
  Target figure:
  \texttt{results/fig\_Hperp\_denominator\_controls\_with\_124\_N3\_N40\_stage123.png}
  (8×5 grid; the \textbf{vertical axis of each panel is the dormant
  fraction H⊥/H} {[}Eq. 23; meaning in §2.5{]} on a semilog scale,
  horizontal axis is the step; the six denominator curves are overlaid.
  Straight segments on the semilog correspond to inflation-like
  exponential amplification; horizontal segments to saturation.)
\item
  The complex-plane readout figures (step 0, final, zoom) are treated in
  Chapter 3.
\end{itemize}

\subsection{6. Execution Analysis (objective report and observations
only; numbers
from}\label{execution-analysis-objective-report-and-observations-only-numbers-from}

analysis\_sweep\_summary\_v1.json)

\begin{enumerate}
\def\labelenumi{\arabic{enumi}.}
\tightlist
\item
  Of the 228 runs, \textbf{212 crossed H⊥/H \textgreater{} 0.05 within
  500 steps} (onset range 45--481).
\item
  \textbf{In all 228 runs, step 1 remained at seed scale}: H⊥/H(0) ∈
  {[}1.92e-33, 1.60e-31{]}, H⊥/H(1) ∈ {[}3.64e-30, 7.32e-25{]}. There
  were \textbf{zero} mismatch injections (jumps of order
  10\textsuperscript{-8--10}-3).
\item
  Final values of the crossing runs: 0.0501--0.458.
\item
  Breakdown of the 16 non-crossing runs: 15 runs saturated with final
  0.025--0.048, \textbf{just below the 0.05 threshold} (the 124 series
  and N+2 series for N≥26, and the three series of N=39); only one run
  was actually slow-growing: N=3, den=124 (final 1.75e-8).
\item
  The onset of the 124 series shortens from 318 at N=4 as N increases,
  entering the band 89--160 for N≥15 (crossing N only; the table is
  \texttt{onset\_by\_series{[}\textquotesingle{}124\textquotesingle{}{]}}
  of \texttt{analysis\_sweep\_summary\_v1.json}).
\item
  The zero closure \textbar z\^{}Tz\textbar/H stayed at most 3.72e-13
  over all 228×501 recorded points (initial max 5.19e-15) --- a
  numerical confirmation of the conservation law of Eq. 19.
\item
  Visual observation of the figure: every panel shows the same family of
  curves --- ``straight exponential amplification from seed scale →
  saturation in the 10\textsuperscript{-3--10}-1 range''.
\end{enumerate}

\subsection{Appendix A: Provenance of the Stage Composition --- Audit
Table of the
Single-Factor}\label{appendix-a-provenance-of-the-stage-composition-audit-table-of-the-single-factor}

Experiments and Deletion Controls

The dynamics of this chapter (stages 1+2+3) is the composition
established by the single-factor experiment series on N=40 with the
static parent (the same Z0 as Chapter 1). The complete set of
experiments (programs, data, figures, README, SHA256SUMS) is contained
in \texttt{ChatGPT\_denominator\_controls\_N40\_selfcontrol\_20260904/}
and is \textbf{bundled with the upload of this paper}. The numbers
originate from the summary CSVs and READMEs of the respective results
directories.

{\def\LTcaptype{none} % do not increment counter
\begin{longtable}[]{@{}
  >{\raggedright\arraybackslash}p{(\linewidth - 10\tabcolsep) * \real{0.1667}}
  >{\raggedright\arraybackslash}p{(\linewidth - 10\tabcolsep) * \real{0.1667}}
  >{\raggedright\arraybackslash}p{(\linewidth - 10\tabcolsep) * \real{0.1667}}
  >{\raggedright\arraybackslash}p{(\linewidth - 10\tabcolsep) * \real{0.1667}}
  >{\raggedright\arraybackslash}p{(\linewidth - 10\tabcolsep) * \real{0.1667}}
  >{\raggedright\arraybackslash}p{(\linewidth - 10\tabcolsep) * \real{0.1667}}@{}}
\toprule\noalign{}
\begin{minipage}[b]{\linewidth}\raggedright
Composition
\end{minipage} & \begin{minipage}[b]{\linewidth}\raggedright
Dynamics (generator / clock)
\end{minipage} & \begin{minipage}[b]{\linewidth}\raggedright
Level of f(1)
\end{minipage} & \begin{minipage}[b]{\linewidth}\raggedright
0.05 crossing (500 steps, 6 dens)
\end{minipage} & \begin{minipage}[b]{\linewidth}\raggedright
Final values / observation
\end{minipage} & \begin{minipage}[b]{\linewidth}\raggedright
Output directory
\end{minipage} \\
\midrule\noalign{}
\endhead
\bottomrule\noalign{}
\endlastfoot
Stage 1 only (baseline) & H=A*(z̄⊗z) (amplitude-weighted, cos included) /
fixed Δτ & 5.3--6.5e-8 (injection) & 0/6 & stable at ceiling
\textasciitilde1.16e-3 & results\_staticparent/ \\
Stages 1+2 (stage 3 deleted) & phase-only Ĥ (cos included) / fixed Δτ &
1.4e-9--8.9e-3 (injection, den-dependent) & 6/6 (τ=4--198) & departure
up to 0.94--0.9999 & results\_staticparent\_phaseonly/ \\
Stages 1+3 (stage 2 deleted) & amplitude-weighted i·Im(H) / fixed Δτ &
1.8--2.3e-8 (injection) & 0/6 & slow drift to \textasciitilde3.5e-3 &
results\_staticparent\_ampimK/ \\
\textbf{Stages 1+2+3 (this chapter)} & phase-only i·K / fixed Δτ &
\textbf{4.0--9.1e-29 (seed scale)} & 4/6 (τ=276--456) & 0.038--0.10,
relaxation curve present & results\_staticparent\_imK/ \\
Stage 2 at init only (reference) & z0 equimodularized once +
amplitude-weighted i·Im(H) / fixed Δτ & 2.9--3.4e-3 (injection) & 6/6
(τ=3--6) & 0.96--0.99, immediate departure &
results\_staticparent\_stage2init/ \\
Stages 2+3 + σ clock (reference) & phase-only i·K / σ-normalized Cayley
clock & 5.84e-30 (seed scale) & 0/6 (within 500 steps) & late-segment
slope 65.8 steps/decade (July canonical: 64.0) &
results\_staticparent\_sigmaclock/ \\
\end{longtable}
}

Observation (facts only): f(1) remains at seed scale
(\textasciitilde10\^{}-29) \textbf{only} in the composition where stages
2 and 3 are present \textbf{simultaneously}; deleting either one makes
f(1) jump to 10\textsuperscript{-8--10}-3. Moving stage 2 to a one-time
initialization produces the largest jump (10\^{}-3).

\subsection{Appendix B: Program Lineage (chain of minimal diffs and
SHA256)}\label{appendix-b-program-lineage-chain-of-minimal-diffs-and-sha256}

The diffs of each step are recorded in full in the READMEs of the
respective packages.

{\def\LTcaptype{none} % do not increment counter
\begin{longtable}[]{@{}
  >{\raggedright\arraybackslash}p{(\linewidth - 4\tabcolsep) * \real{0.3333}}
  >{\raggedright\arraybackslash}p{(\linewidth - 4\tabcolsep) * \real{0.3333}}
  >{\raggedright\arraybackslash}p{(\linewidth - 4\tabcolsep) * \real{0.3333}}@{}}
\toprule\noalign{}
\begin{minipage}[b]{\linewidth}\raggedright
\#
\end{minipage} & \begin{minipage}[b]{\linewidth}\raggedright
Program
\end{minipage} & \begin{minipage}[b]{\linewidth}\raggedright
Diff from the previous
\end{minipage} \\
\midrule\noalign{}
\endhead
\bottomrule\noalign{}
\endlastfoot
0 &
\texttt{ChatGPT\_denominator\_controls\_N3\_N40\_mixedseed\_20260903/run\_and\_plot\_N3\_N40\_mixedseed\_20260903.py}
(ChatGPT-authored canonical) & --- \\
1 &
\texttt{ChatGPT\_denominator\_controls\_N40\_selfcontrol\_20260904/run\_and\_plot\_N40\_only\_selfcontrol.py}
& 2 lines: OUT destination, \texttt{range(40,41)} (N=40 output verified
bit-identical to canonical) \\
2 & \texttt{…/run\_N40\_staticparent\_v1.py} & initial data replaced by
the Chapter-1 static parent Z0 (input only) \\
3 & \texttt{…/run\_N40\_staticparent\_phaseonly\_v1.py} & 1 line in
one\_step: H built from exp(i·arg z) (stage 2) \\
4 & \texttt{…/run\_N40\_staticparent\_imK\_v1.py} & 1 line in one\_step:
H=1j·Im(Ĥ) (stage 3) \\
5 &
\texttt{N3\_N40\_stage123\_sweep\_20260905/run\_N3\_N40\_stage123\_v1.py}
(this chapter) & loop range(3,41), per-N parent loading, output
destination, figure/metadata names only (dynamics unchanged) \\
\end{longtable}
}

SHA256 (full digits; transcribed from the output of shasum -a 256):

\begin{verbatim}
c709c56335d4c67373bff9a3ef6414ea17564d5fbf9c10b7bc9c3724ff091b92  run_and_plot_N3_N40_mixedseed_20260903.py
5a07f354e19985dca0f5de89217e2aa22ac511afaa4b2bb4aa5d93e0c7f9706f  run_and_plot_N40_only_selfcontrol.py
cb5a0ab6db9ae5719eac7aee3b539eb04ca09f5ebfcab6dcce36e8a6e727719e  run_N40_staticparent_v1.py
c1d6d2e60e101f0a99585eefdc22990a087c46246aba112042ac97a7fcbd1a71  run_N40_staticparent_phaseonly_v1.py
a67912b77f7f112731c1eac7612f21b464aed7e2c42431f7b3a7afabb2cd051d  run_N40_staticparent_imK_v1.py
1abf2353fee2e4f56f05e7a6f149fd086885136beb61ab571b48a56b09691567  run_N3_N40_stage123_v1.py
\end{verbatim}

\begin{center}\rule{0.5\linewidth}{0.5pt}\end{center}

(End of Chapter 2. Chapter 3, ``Complex-Plane Readout Figures'',
follows.)

\end{document}
